\chapter{Planning, Motion Planning, and TAMP}

\section*{학습 목표}

이 장은 Real VLA에서 planning layer가 왜 사라지지 않는지 설명한다. VLA가 action을 직접 생성하더라도 long-horizon task에서는 goal decomposition, precondition, feasibility, recovery, constraint satisfaction이 필요하다. 이 장의 목표는 classical planning을 복고적으로 소개하는 것이 아니라, VLA가 실제 로봇 시스템 안에서 planner와 어떻게 결합되는지 이해하는 것이다.

\begin{enumerate}
  \item symbolic planning, motion planning, TAMP의 문제 정의를 구분한다.
  \item language instruction이 plan, subgoal, constraint로 바뀌는 과정을 정식화한다.
  \item continuous feasibility가 왜 LLM/VLA planning의 병목인지 설명한다.
  \item VLA planner, affordance model, motion planner, controller가 하나의 loop로 연결되는 구조를 정리한다.
\end{enumerate}

\section{Planning은 action sequence의 문제다}

자연어 명령은 종종 하나의 문장으로 주어진다. 그러나 로봇이 실행해야 하는 것은 보통 하나의 action이 아니라 조건을 만족하는 action sequence이다. 예를 들어 ``테이블 위 컵을 식기건조대에 놓아라''라는 명령은 object grounding, grasp feasibility, collision-free transport, stable placement, release verification을 포함한다. 이 문제를 바로 continuous action vector 하나로 바꾸면 중요한 구조가 사라진다.

task planning 문제는 다음 tuple로 표현할 수 있다 \cite{lavalle2006,kaelbling2011tamp}.
\begin{equation}
  \mathcal{P}
  =
  (\mathcal{S},\mathcal{A},s_0,\mathcal{G},\gamma).
  \label{eq:planning-problem}
\end{equation}
여기서 $\mathcal{S}$는 symbolic state set, $\mathcal{A}$는 symbolic action set, $s_0$는 initial state, $\mathcal{G}$는 goal condition, $\gamma$는 transition function이다. 하나의 plan은 action sequence이다.
\begin{equation}
  \pi^{plan}
  =
  (a_1,a_2,\ldots,a_K),
  \qquad
  s_K \in \mathcal{G}.
  \label{eq:plan-sequence}
\end{equation}
이때 $a_k$는 VLA policy의 continuous action $\act_t$와 구분하기 위해 symbolic action으로 읽어야 한다. 예를 들어 \texttt{pick(mug)}, \texttt{place(mug,rack)}, \texttt{open(drawer)}는 symbolic action이고, end-effector delta pose나 gripper command는 controller가 실행하는 action이다.

\begin{definitionbox}{Task planning}
Task planning은 목표를 달성하기 위한 discrete action sequence를 찾는 문제이다. 이 계층은 무엇을 어떤 순서로 해야 하는지 다루지만, 각 action이 실제 로봇의 configuration space에서 가능한지는 별도 검사가 필요하다.
\end{definitionbox}

\section{Preconditions and effects}

고전 symbolic planning의 핵심은 action을 precondition과 effect로 표현하는 것이다. action $a$가 state $s$에서 실행 가능하려면 precondition이 만족되어야 한다.
\begin{equation}
  \mathrm{Pre}(a) \subseteq s.
  \label{eq:precondition}
\end{equation}
action을 실행한 뒤의 state는 add effect와 delete effect로 갱신된다.
\begin{equation}
  \gamma(s,a)
  =
  (s \setminus \mathrm{Del}(a)) \cup \mathrm{Add}(a).
  \label{eq:add-delete-effect}
\end{equation}

로봇 조작 예제에서 \texttt{pick(mug)}의 precondition은 다음과 같은 symbolic predicate로 구성될 수 있다.
\begin{equation}
  \begin{aligned}
  \mathrm{Pre}(\texttt{pick(mug)})
  =
  \{&
  \texttt{visible(mug)},\;
  \texttt{reachable(mug)},\\
  &
  \texttt{graspable(mug)},\;
  \texttt{gripper-empty}
  \}.
  \end{aligned}
  \label{eq:pick-preconditions}
\end{equation}

문제는 \texttt{reachable}이나 \texttt{graspable}이 순수 symbolic predicate가 아니라는 점이다. 이것들은 robot geometry, scene geometry, grasp sampler, IK, collision checking, contact condition에 의해 결정된다. 이 지점에서 task planning과 motion planning이 결합된다.

\section{Search and heuristic}

Planning은 search problem으로 볼 수 있다. 각 node는 symbolic state이고, edge는 action이다. 비용이 있는 planning에서는 다음 목적을 최소화한다.
\begin{equation}
  \pi^{plan\star}
  =
  \arg\min_{\pi^{plan}}
  \sum_{k=1}^{K} c(s_{k-1},a_k)
  \quad
  \text{s.t.}\quad
  s_K\in\mathcal{G}.
  \label{eq:planning-cost}
\end{equation}
A*와 같은 heuristic search는 다음 priority를 사용한다.
\begin{equation}
  F(n)=G(n)+H(n),
  \label{eq:a-star}
\end{equation}
여기서 $G(n)$은 현재 node까지의 cost, $H(n)$은 goal까지의 추정 cost이다.

LLM/VLM planner를 이 관점으로 보면, 큰 모델은 heuristic 또는 proposal generator 역할을 할 수 있다. 즉, 어떤 action sequence가 plausible한지 제안할 수 있다. 그러나 plausible plan과 executable plan은 다르다. Real VLA에서는 language model proposal이 feasibility checker, affordance model, motion planner를 통과해야 한다.

\begin{warningbox}{Plausible plan은 executable plan이 아니다}
LLM이 ``컵을 집고 선반에 놓아라''라는 plan을 생성할 수 있어도, 컵이 reachable한지, grasp가 가능한지, 선반까지 collision-free path가 있는지, 놓은 뒤 안정적인지 보장하지 않는다. 자연어 계획은 physical feasibility check를 대체하지 못한다.
\end{warningbox}

\section{Motion planning in configuration space}

Motion planning은 continuous configuration space에서 충돌 없는 경로를 찾는 문제이다. 로봇 configuration을 $\conf\in\mathcal{Q}$라 하자. 장애물과 self-collision을 피하는 자유공간은
\begin{equation}
  \mathcal{Q}_{free}
  =
  \{\conf\in\mathcal{Q}\mid \mathrm{Collision}(\conf)=0\}.
  \label{eq:q-free}
\end{equation}
motion planning 문제는 시작 configuration $\conf_{start}$와 목표 set $\mathcal{Q}_{goal}$ 사이의 경로를 찾는 것이다.
\begin{equation}
  \xi:[0,1]\rightarrow \mathcal{Q}_{free},
  \qquad
  \xi(0)=\conf_{start},
  \qquad
  \xi(1)\in\mathcal{Q}_{goal}.
  \label{eq:motion-planning-path}
\end{equation}

VLA가 end-effector target이나 waypoint를 출력한다면, 그것은 보통 $\mathcal{Q}_{goal}$ 또는 intermediate constraint를 정의한다. 하지만 실제 path $\xi$는 collision, joint limit, velocity limit, singularity, workspace bound를 만족해야 한다.

\section{Sampling-based planning}

PRM이나 RRT 계열의 sampling-based planner는 continuous space를 직접 격자로 나누지 않고 sampled configuration을 이용해 graph 또는 tree를 만든다. 단순화하면 roadmap planner는 다음 graph를 만든다.
\begin{equation}
  \begin{aligned}
  \mathcal{G}_{roadmap}&=(V,E),\qquad
  V\subset\mathcal{Q}_{free},\\
  E&=
  \{(q_i,q_j)\mid
  \mathrm{LocalPlan}(q_i,q_j)\subset\mathcal{Q}_{free}\}.
  \end{aligned}
  \label{eq:roadmap}
\end{equation}

RRT 계열은 시작 configuration에서 tree를 확장한다.
\begin{equation}
  q_{new}
  =
  \mathrm{Steer}(q_{near},q_{rand},\eta),
  \qquad
  q_{new}\in\mathcal{Q}_{free}.
  \label{eq:rrt-steer}
\end{equation}

이 방법들은 high-dimensional configuration space에서 유용하지만, VLA와 결합할 때 두 가지 문제가 생긴다. 첫째, planner가 실패했을 때 그것이 VLA의 잘못인지, goal pose의 잘못인지, scene estimate의 잘못인지 구분해야 한다. 둘째, planner의 computation time이 VLA inference loop와 맞아야 한다. 따라서 planner output은 trajectory뿐 아니라 failure reason을 함께 반환해야 한다.

\section{Trajectory optimization}

Trajectory optimization은 경로를 변수로 두고 objective와 constraint를 최적화한다.
\begin{equation}
\begin{aligned}
  \xi^\star
  =
  \arg\min_{\xi}
  &\quad
  J_{task}(\xi;\act_t)
  +\lambda_s J_{smooth}(\xi)
  +\lambda_c J_{collision}(\xi) \\
  \text{s.t.}
  &\quad
  \xi(\tau)\in\mathcal{Q}_{free},\\
  &\quad
  \dot{\xi}(\tau)\in\mathcal{V},\quad
  \ddot{\xi}(\tau)\in\mathcal{A}.
\end{aligned}
\label{eq:trajopt}
\end{equation}

Trajectory optimization은 VLA가 생성한 goal, cost map, affordance map, forbidden region을 objective나 constraint로 넣기 좋다. 예를 들어 ``유리컵을 건드리지 말라''는 instruction은 forbidden region $\mathcal{R}_{forbid}$로 바뀔 수 있고, collision 또는 distance constraint로 들어갈 수 있다.
\begin{equation}
  d(\xi(\tau),\mathcal{R}_{forbid}) \ge d_{min}
  \quad
  \forall \tau.
  \label{eq:forbidden-distance}
\end{equation}

\section{Task and motion planning}

TAMP는 symbolic action sequence와 continuous parameter를 함께 찾는 문제이다. symbolic action $a_k$는 grasp, placement, path, timing 같은 continuous parameter $\theta_k$를 필요로 한다.
\begin{equation}
  (a_{1:K}^{\star},\theta_{1:K}^{\star})
  =
  \arg\min_{a_{1:K},\theta_{1:K}}
  \sum_{k=1}^{K}
  C(a_k,\theta_k)
  \quad
  \text{s.t.}\quad
  \Phi(a_{1:K},\theta_{1:K};\hat{\state}_0)=1.
  \label{eq:tamp}
\end{equation}
여기서 $\Phi$는 symbolic precondition, geometric feasibility, collision checking, grasp validity, placement stability를 모두 포함하는 feasibility predicate이다.

예를 들어 \texttt{pick(mug)}는 단순한 symbolic action처럼 보이지만 실제로는 grasp pose $g$, pregrasp pose $p$, approach trajectory $\xi$, gripper force $f_g$를 필요로 한다.
\begin{equation}
  \theta_{\texttt{pick}}
  =
  (g,p,\xi,f_g).
  \label{eq:pick-continuous-params}
\end{equation}
따라서 TAMP는 ``무엇을 할 것인가''와 ``그것을 어떻게 물리적으로 실행할 것인가''를 동시에 다룬다.

\begin{table}[ht]
\centering
\caption{Planning 계층별 책임}
\label{tab:planning-layers}
\small
\begin{tabularx}{\textwidth}{p{30mm}X X}
\toprule
계층 & 입력/출력 & 대표 failure \\
\midrule
Task planning & symbolic state, goal $\rightarrow$ action sequence & missing precondition, wrong ordering, no recovery branch \\
Motion planning & start/goal configuration, geometry $\rightarrow$ collision-free path & IK failure, narrow passage, collision, timeout \\
Trajectory optimization & cost, constraints $\rightarrow$ smooth feasible trajectory & local minimum, constraint violation, solver latency \\
TAMP & symbolic action + continuous parameter search & sampled grasp/pose infeasible, combinatorial explosion \\
VLA/ER planning & language, image, memory $\rightarrow$ subgoal, affordance, plan proposal & plausible but infeasible plan, hallucinated object, unsafe ambiguity \\
\bottomrule
\end{tabularx}
\end{table}

\section{VLA as planner, scorer, or executor}

VLA와 planning의 결합 방식은 하나가 아니다. Real VLA stack에서 큰 모델은 다음 네 역할 중 하나를 맡을 수 있다.

\begin{enumerate}
  \item \textbf{Planner}: language와 scene을 보고 symbolic 또는 semi-symbolic plan을 제안한다.
  \item \textbf{Scorer}: 여러 skill 또는 action 후보의 affordance score를 계산한다.
  \item \textbf{Constraint generator}: forbidden region, target relation, success condition을 만든다.
  \item \textbf{Executor}: observation과 subgoal을 받아 continuous action을 생성한다.
\end{enumerate}

이 역할들은 한 시스템 안에서 함께 존재할 수 있다.
\begin{equation}
  \task,\obs_t
  \xrightarrow{\text{ER/VLM}}
  (g_{1:K},\mathcal{C}_{1:K})
  \xrightarrow{\text{TAMP/Motion Planner}}
  (\xi_{1:K},r_{1:K})
  \xrightarrow{\text{VLA/Controller}}
  \ctrlcmd_t.
  \label{eq:vla-planning-pipeline}
\end{equation}
여기서 $g_{1:K}$는 subgoal sequence, $\mathcal{C}_{1:K}$는 constraint sequence, $\xi_{1:K}$는 feasible trajectory, $r_{1:K}$는 controller reference이다.

\section{Closed-loop replanning}

실제 로봇 실행은 open-loop plan execution이 아니다. perception error, object motion, human intervention, contact failure, planner timeout이 발생하면 replanning이 필요하다. 하나의 receding planning loop는 다음과 같이 쓸 수 있다.
\begin{equation}
\begin{aligned}
  \bel_t &= \mathrm{UpdateBelief}(\bel_{t-1},\obs_t),\\
  p_t &= \mathrm{Plan}(\task,\bel_t,\mathcal{C}_t),\\
  \act_t &= \policy(\obs_t,\conf_t,p_t),\\
  \ctrlcmd_t &= \controller(\act_t,\hat{\state}_t,\mathcal{C}_t).
\end{aligned}
\label{eq:closed-loop-planning}
\end{equation}

이 loop에서 VLA는 plan을 매 step 다시 만들 수도 있고, plan은 느리게 갱신하고 VLA policy만 빠르게 실행할 수도 있다. 핵심은 plan, action, controller, safety가 서로 다른 rate로 동작할 수 있다는 점이다.

\section{Evaluation metadata}

Planning이 포함된 VLA 실험은 단순 success rate만으로 평가하면 부족하다. 최소한 다음 항목을 기록해야 한다.

\begin{itemize}
  \item plan length and number of replans
  \item symbolic failure reason: missing object, unsatisfied precondition, no valid action
  \item motion planning failure reason: IK failure, collision, timeout, narrow passage
  \item trajectory optimization status: converged, infeasible, constraint violation
  \item VLA role: planner, scorer, constraint generator, executor
  \item recovery branch and human intervention count
  \item time budget: ER/VLM inference, planning, motion planning, controller loop
\end{itemize}

이 기록이 없으면 ``VLA가 실패했다''라는 문장은 너무 거칠다. 실패는 semantic parsing, object grounding, symbolic precondition, grasp sampling, IK, collision checking, controller tracking, contact safety 중 어느 곳에서든 발생할 수 있다.

\begin{casebox}{Case study: SayCan과 VoxPoser를 feasibility 관점에서 읽기}
SayCan은 language likelihood와 affordance score를 결합해 executable skill을 고르려 했고, VoxPoser는 language-conditioned value map을 통해 motion optimization이 사용할 field를 만들었다 \cite{saycan2022,voxposer2023}. 두 사례 모두 핵심은 planner가 자연어 plan을 쓰는지가 아니라, semantic proposal이 feasibility check와 어떻게 만나는가이다.
\end{casebox}

\chapterwork
{Sampling-based motion planning과 TAMP를 LLM/VLM planner 논문과 함께 읽어라 \cite{lavalle2006,kavraki1996,lavalle2001rrt,kaelbling2011tamp}.}
{LLM이 plan text를 잘 생성해도 robot execution plan이 되지 않는 이유는 무엇인가?}
{mug/rack/glass task를 symbolic precondition, motion feasibility, contact constraint, recovery branch로 나누어 plan schema를 작성하라.}

\section{요약}

Planning은 Real VLA에서 과거의 모듈이 아니라 현재의 responsibility layer이다. LLM/VLM/VLA는 plan proposal, affordance scoring, constraint generation, action execution을 강하게 만들 수 있다. 그러나 physical feasibility, collision checking, contact constraint, recovery planning은 여전히 시스템 안에 명시적으로 존재해야 한다. 강한 Real VLA는 end-to-end policy와 planner 중 하나를 선택하는 시스템이 아니라, semantic reasoning과 geometric feasibility를 closed-loop로 연결하는 시스템이다.

다음 장에서는 policy를 어떻게 학습하는가로 이동한다. planning이 무엇을 해야 하는지와 어떤 경로가 가능한지를 다룬다면, imitation learning은 실제 trajectory에서 action policy를 어떻게 배우는지 다룬다.
